% Author: Prof. Dr. Matthias Jung, DL9MJ
% Year: 2026
%
% Strahlungsmuster und Stromverteilung idealer Monopolantennen
% über perfekter Erde.
%
% Elevationswinkel alpha:
%   0°  = Horizont
%   90° = Zenit
%
% Gezeigt werden:
%   lambda/4
%   lambda/2
%   5 lambda/8
%   lambda
%
% Die orangen Kurven zeigen die relative elektrische Feldstärke |E|
% im Fernfeld bei gleicher insgesamt abgestrahlter Leistung.
%
% Alle Diagramme sind gemeinsam auf den lambda/4-Monopol normiert:
%   E_max(lambda/4) = 1
%
% Die roten Kurven zeigen die idealisierte Stromverteilung
% entlang des Strahlers.
%
% Für die Feldstärke ergeben sich ungefähr:
%
%   lambda/4:      E_max = 1.000
%   lambda/2:      E_max = 1.212
%   5 lambda/8:    E_max = 1.414
%   lambda:        E_max = 1.241
%

\begin{tikzpicture}

% --------------------------------------------------------------------------
% Fernfeld eines idealen vertikalen Monopols
%
% alpha = Elevationswinkel über dem Horizont
% h     = Strahlerlänge in Wellenlängen
%
% Unnormierter Feldfaktor:
%
%                  cos(2*pi*h*sin(alpha)) - cos(2*pi*h)
% F(alpha) = ---------------------------------------------------
%                              cos(alpha)
%
% Die Antennen unterschiedlicher Länge werden so normiert,
% dass jeweils dieselbe Gesamtleistung abgestrahlt wird.
%
% Als gemeinsame Referenz dient der lambda/4-Monopol.
%
% Die Normierungsfaktoren für die Feldstärke sind die
% Quadratwurzeln der entsprechenden Leistungsfaktoren:
%
%   lambda/4:      1.000000000
%   lambda/2:      0.606072220
%   5 lambda/8:    0.828507649
%   lambda:        0.530720330
%
% PGF verwendet Grad bei trigonometrischen Funktionen.
% Daher entspricht 2*pi*h dem Winkel 360*h.
% --------------------------------------------------------------------------

\pgfmathdeclarefunction{monofield}{3}{%
    \pgfmathparse{%
        abs(
            (
                cos(360*#2*sin(#1))
                - cos(360*#2)
            )
            / cos(#1)
        )
        * #3
    }%
}

% --------------------------------------------------------------------------
% Polardiagramm
%
% Eine radiale Einheit entspricht dem maximalen E-Feld
% des lambda/4-Monopols.
%
% Der äußere Radius entspricht E = 1.5.
% --------------------------------------------------------------------------

\def\R{0.93}
\def\Emax{1.5}

% Horizontale Auslenkung der Stromkurve
\def\Iscale{0.28}

% Vertikale Skalierung der Antennendarstellung
\def\AntennaScale{1.55}

\tikzset{
    polargrid/.style={
        draw=black!35,
        line width=0.35pt
    },
    polargridminor/.style={
        draw=black!20,
        line width=0.25pt
    },
    polarreference/.style={
        draw=black!45,
        line width=0.45pt
    },
    radiation/.style={
        draw=DARCorange,
        line width=1.4pt
    },
    antenna/.style={
        draw=black,
        line width=1.1pt
    },
    current/.style={
        draw=DARCred,
        line width=1.3pt
    }
}

% --------------------------------------------------------------------------
% Panel
%
% #1 x-Position
% #2 y-Position
% #3 Strahlerlänge h/lambda
% #4 Normierungsfaktor des E-Feldes
% #5 Beschriftung der mechanischen Länge
% --------------------------------------------------------------------------

\newcommand{\MonopolePanel}[5]{%
    \begin{scope}[shift={(#1,#2)}]

        % --------------------------------------------------------------
        % Polardiagramm
        %
        % Radiale Kreise:
        %   0.5
        %   1.0   = Maximum lambda/4
        %   1.5
        % --------------------------------------------------------------

        \foreach \r in {0.5,1.5}{
            \draw[polargridminor]
                (0,0)
                ++(0:{\r*\R})
                arc[
                    start angle=0,
                    end angle=90,
                    radius={\r*\R}
                ];
        }

        % Referenzkreis E = 1
        \draw[polarreference]
            (0,0)
            ++(0:\R)
            arc[
                start angle=0,
                end angle=90,
                radius=\R
            ];

        % Winkelstrahlen
        \foreach \a in {0,10,...,90}{
            \draw[polargrid]
                (0,0) -- (\a:{\Emax*\R});
        }

        % Winkelbeschriftungen
        \foreach \a in {0,10,...,90}{
            \node[
                font=\scriptsize,
                inner sep=1pt
            ] at (\a:{1.60*\R}) {\a};
        }

        % --------------------------------------------------------------
        % Strahlungsdiagramm
        %
        % Radius proportional zur relativen elektrischen Feldstärke |E|.
        %
        % Alle vier Diagramme benutzen dieselbe Referenz:
        %
        %       E_max(lambda/4) = 1
        % --------------------------------------------------------------

        \draw[radiation]
            plot[
                domain=0:89.7,
                samples=350,
                smooth,
                variable=\a
            ]
            (
                {
                    \R
                    * monofield(\a,#3,#4)
                    * cos(\a)
                },
                {
                    \R
                    * monofield(\a,#3,#4)
                    * sin(\a)
                }
            );

        % --------------------------------------------------------------
        % Antennendarstellung
        % --------------------------------------------------------------

        \pgfmathsetmacro{\antennaheight}{\AntennaScale*#3}

        % Groundplane / Erde
        \draw[antenna]
            (-0.73,0) -- (-0.27,0);

        % Strahler
        \draw[antenna]
            (-0.50,0) -- (-0.50,\antennaheight);

        % --------------------------------------------------------------
        % Stromverteilung
        %
        % I(z) = Imax * sin(k * (l-z))
        %
        % Da z und l in Wellenlängen angegeben sind:
        %
        % I(z) = sin(360° * (l-z))
        %
        % Die Stromverteilung ist für jede Antenne auf denselben
        % zeichnerischen Ausschlag skaliert. Die roten Kurven zeigen
        % daher Form und Phasenlage, nicht die absolute Stromstärke.
        %
        % Ein Vorzeichenwechsel erscheint auf der jeweils anderen
        % Seite des Strahlers.
        % --------------------------------------------------------------

        \draw[current]
            plot[
                domain=0:#3,
                samples=180,
                smooth,
                variable=\z
            ]
            (
                {
                    -0.50
                    + \Iscale*sin(360*(#3-\z))
                },
                {
                    \AntennaScale*\z
                }
            );

        % --------------------------------------------------------------
        % Längenmaß
        % --------------------------------------------------------------

        \draw[
            Triangle-Triangle,
            line width=0.55pt
        ]
            (-0.82,0)
            --
            (-0.82,\antennaheight);

        % Längenbeschriftung
        \node[
            rotate=90,
            font=\small,
            anchor=south
        ] at (-0.87,{0.5*\antennaheight})
            {$l=#5$};

    \end{scope}
}

% --------------------------------------------------------------------------
% Obere Reihe
% --------------------------------------------------------------------------

% lambda/4
%
% E_max = 1.000 bei alpha = 0°
%
\MonopolePanel
    {0}
    {0}
    {0.25}
    {1.000000000}
    {\frac{1}{4}\lambda}


% lambda/2
%
% E_max ~= 1.212 bei alpha = 0°
%
\MonopolePanel
    {3.15}
    {0}
    {0.50}
    {0.606072220}
    {\frac{1}{2}\lambda}


% --------------------------------------------------------------------------
% Untere Reihe
% --------------------------------------------------------------------------

% 5 lambda/8
%
% E_max ~= 1.414 bei alpha = 0°
%
% erste Nullstelle:
% alpha ~= 36.87°
%
% kleine Hochwinkelkeule
%
\MonopolePanel
    {0}
    {-2.05}
    {0.625}
    {0.828507649}
    {\frac{5}{8}\lambda}


% lambda
%
% Nullstelle am Horizont:
% alpha = 0°
%
% E_max ~= 1.241
% Hauptmaximum bei alpha ~= 32.56°
%
\MonopolePanel
    {3.15}
    {-2.05}
    {1.0}
    {0.530720330}
    {\lambda}

\end{tikzpicture}